\documentclass[11pt]{article}
\usepackage[margin=1in]{geometry}
\usepackage{amsmath,amssymb,booktabs,tabularx}
\usepackage[T1]{fontenc}
\usepackage{lmodern}
\usepackage{microtype}
\usepackage[hidelinks]{hyperref}
\usepackage{enumitem}
\setlist{nosep}
\setlength{\parindent}{0pt}
\setlength{\parskip}{6pt}
\title{Dynamic NBA Team Ratings and Pregame Win Probabilities\\
\large Model Specification and Evaluation Protocol}
\author{Eshan Garkhedkar}
\date{Working draft 0.1 --- October 12, 2026}

\begin{document}
\maketitle
\begin{abstract}
This project develops an NBA forecasting system that maintains team strength
ratings and estimates pregame win probabilities. It adapts a multifactor
Elo-style and Bradley--Terry framework developed for the 2025 Wharton High
School Data Science Competition. The proposed system combines rating
differences with home-court, rest, travel, and time-zone effects, and adjusts
ratings after completed games using schedule strength and performance
information. This draft specifies the model family, data requirements, and
chronological evaluation protocol. NBA parameters have not yet been fitted,
historical data coverage remains unverified, and no NBA predictive performance
is claimed. Equations marked provisional require implementation and validation
before becoming the deployed specification.
\end{abstract}

\section{Scope and provenance}
The intended outputs are a team list ordered by rating, probabilities for
upcoming games, and a permanent record of previously issued predictions and
outcomes. Ratings update after completed games. In-game win probability is
outside the present scope.

The original team report supplies the conceptual framework and feature
definitions \cite{billybeanes}. Associated competition code is reference
material rather than the canonical NBA implementation \cite{referencecode}.
The NBA adaptation will define one consistent probability function, explicit
feature units, reproducible fitting, and timestamped evaluation. Any revision
to the inherited equations will be documented with the model version.

\section{Forecast timing and notation}
For a game $g$ between teams $i$ and $j$, let $\tau_g$ denote its forecast
cutoff and let $\mathcal{F}_{\tau_g}$ contain only information available by
that time. The result is $y_g=1$ when team $i$ wins and $0$ otherwise.
The forecast is
\begin{equation}
 p_g=\Pr(y_g=1\mid\mathcal{F}_{\tau_g}).
\end{equation}
The deployed policy will record the last eligible pregame forecast, its
timestamp, data cutoff, and model version. Historical replays are reported
separately from forecasts genuinely issued before tip-off.

\begin{table}[ht]
\centering\small
\begin{tabularx}{\textwidth}{lX}
\toprule
Symbol & Meaning \\
\midrule
$r_{i,g}$ & Team $i$'s rating immediately before game $g$ \\
$h_{i,g}$ & Venue indicator: $+1$ home, $-1$ away, $0$ neutral \\
$H$ & Home-court offset in rating units \\
$RD_{i,g}$ & Rest days under a fixed calendar convention \\
$TD_{i,g}$ & Estimated distance from the preceding venue, in kilometres \\
$TZ_{i,g}$ & Dated time-zone offset change, in hours \\
$S_{i,g}$ & Opponent-based schedule-strength feature \\
$L_{i,g}$ & Lead-retention feature observed after the game \\
$Q_{i,g}$ & Performance relative to prior history, observed after the game \\
\bottomrule
\end{tabularx}
\caption{Core notation. The subscript $g$ is suppressed where unambiguous.}
\end{table}

\section{Pregame probability model}
The inherited fatigue construction is
\begin{align}
 F_{i,g}&=RD_{j,g}+\exp(|TZ_{i,g}|/3)+\nu TD_{i,g},\\
 G_{i,g}&=RD_{i,g}+\exp(|TZ_{j,g}|/3)+\nu TD_{j,g}.
\end{align}
These correspond to opponent rest/team fatigue and team rest/opponent fatigue.
Their exact functional form is provisional. The factor 3 specifies a time-zone
scale in hours, and $\nu$ converts the travel contribution to the numerical
scale of this index. These are empirical indices, not measurements of
physiological fatigue.

Using natural logarithms, the proposed canonical probability is
\begin{equation}
 p_{i,g}=\left[1+b^{z_{i,g}}\right]^{-1},\qquad
 z_{i,g}=\frac{r_{j,g}-r_{i,g}-Hh_{i,g}}{c}
       +\gamma\log\frac{F_{i,g}}{G_{i,g}},
 \label{eq:probability}
\end{equation}
where $b>1$, $c>0$, and the fatigue indices must be positive.
The opponent probability is computed as $1-p_{i,g}$. Full-precision
probabilities are retained; rounding is only for display.

The probability can equivalently be written
\begin{equation}
 \operatorname{logit}(p_{i,g})=
 a(r_{i,g}-r_{j,g}+Hh_{i,g})-d\log(F_{i,g}/G_{i,g}),
 \quad a=\frac{\log b}{c},\quad d=\gamma\log b.
\end{equation}
This exposes redundant parameterizations. The implementation will anchor the
rating scale and avoid freely fitting interchangeable scale parameters. A
Bayesian optimization algorithm for constants, if used, does not itself make
the rating model a Bayesian posterior model.

\section{Updates after completed games}
Let $s_{i,g}=2y_{i,g}-1$. The inherited update family is
\begin{equation}
\begin{split}
 r_{i,g}^{+}=r_{i,g}+K(y_{i,g}-p_{i,g})
 &(1+s_{i,g}S_{i,g})\\[-2pt]
 {}\times &(1+\phi\log_2(1+L_{i,g}))
 (1+s_{i,g}\lambda(Q_{i,g}-1)).
\end{split}
\label{eq:update}
\end{equation}
Both teams' updates use the same pregame rating state, their own performance
features, and complementary probabilities. Schedule strength uses past data;
the current result, largest lead, and box score become available only after
completion. No current-game box-score feature enters its pregame forecast.

Equation~\ref{eq:update} is a specification of the inherited model family,
not yet a finalized NBA estimator. Multipliers require finite, nonnegative
values to preserve the direction of the outcome residual. Asymmetric updates
need not conserve the sum of ratings. Initialization, inter-season carryover,
rating centring, and multiplier constraints will be fixed before final testing.

\subsection{Schedule strength}
The original feature has the form
\begin{equation}
 S_{i,g}=\frac{2W_{j,g}+\overline{W}_{\mathcal{O}(j),g}}{3},
\end{equation}
where $W_{j,g}$ is the opponent's pregame win fraction and
$\overline{W}_{\mathcal{O}(j),g}$ summarizes that opponent's prior opponents.
Repeated-opponent weighting, early-season smoothing, and season boundaries
require explicit conventions. Full-season win fractions cannot be substituted
into an earlier game's features.

\subsection{Lead retention and performance}
The original lead-retention ratio is
\begin{equation}
 L_{i,g}=\frac{|PTS_{i,g}-PTS_{j,g}|}{LL_{i,g}},
\end{equation}
where $LL$ denotes the team's largest lead. This is undefined for a team
that never led. Its extension for $LL=0$ will be defined and validated before
fitting, rather than silently converting invalid values to small constants.

For eight box-score features $X_{k,i,g}$, the inherited performance index is
\begin{equation}
 U_{i,g}=\sum_{k=1}^{8}\alpha_k X_{k,i,g},\qquad
 A_{i,g}=\frac{U_{i,g}}{1+U_{i,g}},\qquad
 Q_{i,g}=\frac{A_{i,g}}{\overline{A}_{i,<g}}.
\end{equation}
The proposed reference history is the mean of eligible prior games, excluding
the current game. An opening-game rule remains to be specified. Because net
efficiency can be negative, $U/(1+U)$ is not generally bounded and is singular
at $U=-1$. A stable definition or domain restriction must therefore precede
optimization. Negative net efficiency is valid information and must not be
discarded during numeric conversion.

\section{Data construction and quality control}
The proposed raw data comprise game IDs, teams, tip-off times, venues,
statuses, final scores, team box scores, and largest leads. A versioned venue
lookup supplies coordinates and time zones. NBA endpoints accessible through
the open-source \texttt{nba\_api} ecosystem are candidate sources
\cite{nbaapi}; availability and historical completeness are not yet established.
Data acquisition must use free access. No paid feed is assumed.

Canonical tables contain one record per game and one record per team-game.
Count and fraction fields receive distinct names. Validation checks include
unique keys, two opponents per completed game, score agreement, field types,
and missing-value flags. Genuine blowouts are retained. Source responses and
retrieval timestamps support reproducibility and correction handling.

Rest is provisionally the number of intervening calendar days, making a
back-to-back zero. Travel is estimated between consecutive game venues,
not inferred to be the actual team itinerary. Time-zone differences use dated
UTC offsets. Missing inputs require documented treatment using past data only.

\section{Parameter fitting and evaluation protocol}
For parameters $\theta$, replay the training games in chronological order,
predicting before updating. A proposed objective is mean binary log loss:
\begin{equation}
 \mathcal{L}(\theta)=-\frac{1}{N}\sum_{g=1}^{N}
 \left[y_g\log p_g(\theta)+(1-y_g)\log(1-p_g(\theta))\right].
\end{equation}
Every candidate parameter set requires its own rating replay. Numerical
probability clipping, if needed, is fixed and disclosed. Feature scaling,
imputation, constraints, and any regularization are learned or selected
without access to the final test outcomes.

Subject to data coverage, the proposed split is 2021--22 through 2023--24 for
training, 2024--25 for validation, and 2025--26 for final testing. Parameters
and formula choices are frozen before test evaluation. Ratings may update
after each observed test result, matching the intended live process.

Probability evaluation includes log loss, the Brier score
\begin{equation}
 BS=\frac{1}{N}\sum_{g=1}^{N}(p_g-y_g)^2,
\end{equation}
and calibration plots comparing forecast probabilities with observed win
fractions. Report bin counts and uncertainty alongside calibration estimates.
Winner accuracy is secondary, with a stated tie rule for $p=0.5$. Evaluation
must disclose sample counts, excluded games, and the forecast cutoff policy.

\section{Current evidence and pending results}
As a reference check, the original prediction function reproduced all ten
rounded home-team probabilities in the competition report using the repository's
supplied matchup ratings. This verifies that example's probability outputs;
it does not reproduce the historical rating path or establish forecast accuracy.

No NBA training or test results are available for this draft. The results
section will be completed with fitted parameters, coverage statistics,
chronological evaluation, calibration, and subsequently live forecast results.

\section{Deployment and reproducibility}
A scheduled Python process will ingest completed games, validate features,
update rating state, and save forecasts. Streamlit will display those outputs
without refitting on page visits. Failed ingestion preserves the previous
valid snapshot and exposes its age. Corrected results may require replaying
rating state but must not overwrite previously issued forecasts.

Each release will identify the source-code revision, model configuration,
data snapshot or manifest, preprocessing rules, random seed, and evaluation
interval. This manuscript will be revised alongside material changes to
features, equations, fitting, or evaluation. The repository remains private
during development; manuscript draft status does not imply deployment.

\appendix
\section{Box-score feature derivations}
Let $O$ denote the opponent, with all quantities belonging to one completed
game. Percentages are represented as fractions unless explicitly stated.
Possessions are estimated as
\begin{equation}
 \widehat{P}=FGA+0.44FTA-OREB+TOV.
\end{equation}
The following definitions form the proposed NBA feature contract:
\begin{align}
 X_1&=100(PTS-PTS_O)/\widehat{P},\\
 X_2&=(FGM+0.5\,3PM)/FGA,\\
 X_3&=PTS/[2(FGA+0.44FTA)],\\
 X_4&=OREB/(OREB+DREB_O),\\
 X_5&=DREB/(DREB+OREB_O),\\
 X_6&=AST/TOV,\\
 X_7&=TOV/(FGA+0.44FTA+TOV),\\
 X_8&=STL/\widehat{P}_O.
\end{align}
$X_1$ retains the inherited shared-possession approximation for net
efficiency. Possession estimation conventions must be consistent across
teams and seasons. $X_4$ and $X_5$ explicitly use the opposite rebound
category from the opponent. $X_8$ proposes an opponent-possession denominator,
revising the report's expression that omitted offensive rebounds. These
choices must be reconciled and frozen before fitting. Zero denominators are
flagged and handled by a specified rule, not blanket replacement of all zeros.

\section{Travel derivation}
For venue latitude $\varphi$ and longitude $\ell$ in radians, estimated
great-circle distance is
\begin{align}
 a_v&=\sin^2\!\left(\frac{\varphi_2-\varphi_1}{2}\right)
 +\cos\varphi_1\cos\varphi_2
 \sin^2\!\left(\frac{\ell_2-\ell_1}{2}\right),\\
 TD&=2R\arcsin\sqrt{a_v},\qquad R=6371\ \text{km}.
\end{align}
The computation clips $a_v$ to $[0,1]$ for floating-point safety. Neutral
venues and arena changes require the actual game venue rather than a fixed
franchise home location.

\begin{thebibliography}{9}
\bibitem{billybeanes}
A. Bora, A. Lalwani, E. Garkhedkar, and S. Kulkarni.
\emph{Enhanced Basketball Team Ratings and Win Probabilities:
A Multi-Factor Whole-History Rating and Bradley-Terry Approach}.
The Billy Beanes, Wharton High School Data Science Competition report, 2025.
Advisor: C. Hasija. Team-provided report.
\bibitem{referencecode}
A. Bora. \emph{wharton\_game\_prediction\_model}.
\url{https://github.com/aaravbora22/wharton_game_prediction_model}.
Reviewed revision \texttt{a27759d}, October 12, 2026.
\bibitem{nbaapi}
\emph{nba\_api: An API Client Package to Access the APIs of NBA.com}.
\url{https://github.com/swar/nba_api}. Candidate data-access tooling.
\end{thebibliography}
\end{document}
